\scanpage{118}
et un $x'$ au dessus de $x$ dans le normalisé de $S_x$ dans $\overline L$, d'où un groupe de décomposition $N_{x'}$, et un groupe d'inertie $\unsure{L_{x'}|I_{x'}}$, avec
\[
(26)\qquad N_{x'}/I_{x'}\simeq k(x')
\]
(ext. de type fini de $K$, de degré de tr. $n-1$, si $n=\deg\operatorname{tr}L/K$), et
\[
(27)\qquad I_{x'}\simeq\unsure{T_0|T_\infty}(\overline L)\simeq T_\infty(\unsure{\widetilde K|\overline K})\;-
\]
(par Kummer). \deleted{Notons \unsure{bien, on a}} \inserted{Ces constructions} ne dépendent du choix d'aucun \deleted{de} modèle ($U$, ou $X\ldots$). Notons que
\[
(28)\qquad I_{\unsure{x'|x}}\subset\pi_{L/K}^{\overline L}.
\]
En fait, on peut dire, posant
\[
(29)\qquad\left\{\begin{aligned}
\widetilde S_x&=\operatorname{Spec}\widetilde V,\quad\widetilde V\text{ hensélisé (pas strict) de }V\\
\widetilde U_x&=\widetilde S_x\setminus\{x\}\simeq\operatorname{Spec}\widetilde L_x\quad(\text{où }\widetilde L_x\text{ est le corps des \unsure{fractions} de }\widetilde V)
\end{aligned}\right.
\]
que l'on a un hom \inserted{ext} $\mathcal E_{\widetilde U_x}\longrightarrow\mathcal E_L$, \deleted{\unsure{alors}} dont le choix ``effectif'' dépend de
\margin{$I_{x'}$ \unsure{joue} le rôle d'un ``groupe bout''}
\margin{si Car. $k(x)$ est nulle --- \unsure{sinon quotient du seul habituel}\ldots}
\margin{117}
